\subsection{Static and temporal variance must be tested separately}

A finite network can place more input variance in fixed differences between
neurons and less in temporal fluctuations. Testing only the total input
variance $V_a=\Delta_a+q_a$ would miss this change. Let
$\mathcal M_2(h,V,k)=\E[\phi(X_1)\phi(X_2)]$ for two Gaussian variables
with common mean $h$, marginal variance $V$, and covariance $k$. The
conditional Gaussian construction gives
\begin{equation}
 \bar v_a=\mathcal M_2(h_a,V_a,V_a)
          -\mathcal M_2(h_a,V_a,\Delta_a),\qquad
 \left.\frac{\partial\bar v_a}{\partial\Delta_a}\right|_{V_a,h_a}
 =-A_a.
 \label{eq:partitionderivative}
\end{equation}
The second identity follows from Gaussian integration by parts: differentiating
the shared covariance gives $\E[\phi'(X_1)\phi'(X_2)]$, which equals
$\E_z\alpha_a(z)^2=A_a$ at shared covariance $\Delta_a$.
Thus reallocating input variance toward static disorder lowers the temporal
rate variance along this specified diagnostic direction.

At fixed $V,h$, linearizing the static self-consistency functional gives the
susceptibility $(I-W\diag(A))^{-1}$. Its Perron-mode amplification is about
$6.94$ for the reference law. This sensitivity explains why modest changes
in local operating points can matter appreciably for activity fluctuations.
Actual networks also change $V$ and $h$, and the covariance energy changes
through both its source and its response. These derivative identities
therefore diagnose a possible finite-network mechanism; they do not supply
a complete correction to PR.
